\section{定積分と面積}\label{節:積分符号面積}
$f\ge0$ のとき,リーマン和は長方形の面積の和であり,その極限として曲線の下の面積を定める.では,\,$f$ が負になる部分では,定積分は何を表すのだろうか.\,$f(\xi_k)<0$ となる項 $f(\xi_k)\Delta x_k$ は負なので,その部分の寄与は面積に負号を付けたものになる.したがって,定積分をそのまま図形の面積と呼ぶことはできない.

$f$ が正の部分と $x$ 軸の間の面積を $S_+$,負の部分と $x$ 軸の間の面積を $S_-$ とすれば,
\[
 \int_a^b f(x)\dd x=S_+-S_-,\qquad
 \int_a^b|f(x)|\dd x=S_++S_-
\]
である.どちらも同じグラフから求めるが,符号を保つか,すべて正の面積として足すかが異なる.面積を求めるには,\,$f$ が $0$ になる点を調べ,符号が変わる点で区間を分ければよい.

\begin{example}{面積と定積分}
$y=x-1$ の $0\le x\le3$ における定積分と,グラフと $x$ 軸の間の面積を,三角形の面積から求めよ.
\end{example}
\paragraphbf{方針}
$y=x-1$ が $x$ 軸と交わる $x=1$ で区間を分け,2つの三角形の面積を求める.

\paragraphbf{解答}
定積分は $\frac32$,面積は $\frac52$ である.

\paragraphbf{解法}
グラフは $x=1$ で $x$ 軸と交わる.\,$0\le x\le1$ では,\,$x$ 軸の下側に底辺 $1$,高さ $1$ の三角形ができ,その面積は $\frac12$ である.\,$1\le x\le3$ では,上側に底辺 $2$,高さ $2$ の三角形ができ,その面積は $2$ である.したがって,
\[
 \int_0^3(x-1)\dd x=-\frac12+2=\frac32,
 \qquad
 \int_0^3|x-1|\dd x=\frac12+2=\frac52
\]
となる.図\ref{図:定積分符号}の下側の部分は,定積分では引き,面積では足している.

\begin{figure}[htbp]
\centering
\begin{tikzpicture}[x=2.3cm,y=1.5cm,>=Stealth]
 \fill[red!15] (0,0)--(0,-1)--(1,0)--cycle;
 \fill[blue!15] (1,0)--(3,2)--(3,0)--cycle;
 \draw[->] (-0.2,0)--(3.45,0) node[right]{$x$};
 \draw[->] (0,-1.3)--(0,2.5) node[above]{$y$};
 \draw[very thick,domain=0:3.2] plot (\x,{\x-1});
 \draw[dashed] (3,0)--(3,2)--(0,2);
 \node[below] at (1,0) {$1$};
 \node[below] at (3,0) {$3$};
 \node[left] at (0,-1) {$-1$};
 \node[left] at (0,2) {$2$};
 \node[red!70!black] at (0.3,-0.32) {$S_-$};
 \node[blue!70!black] at (2.4,0.48) {$S_+$};
 \node at (2.45,2.25) {$y=x-1$};
\end{tikzpicture}
\caption{定積分は上側の面積から下側の面積を引く.幾何学的な面積は両方を足す.}
\label{図:定積分符号}
\end{figure}


